%% Section 3: IntSeqBERT
%% English manuscript section.
%% -----------------------------------------------------------------------

\section{IntSeqBERT}
\label{sec:model}

\subsection{Problem Formulation}
\label{sec:formulation}

Let $\mathbf{x} = (x_1, x_2, \ldots, x_L)$ with $x_i \in \mathbb{Z}$ and $L \leq 128$ be a
finite prefix of an OEIS sequence. We adopt \emph{masked sequence modelling}: a random subset
of positions is masked and the model is trained to predict each masked value. Specifically, for
each masked position $i$, three quantities are predicted:
\begin{enumerate}
  \item \textbf{Magnitude}: $v_i =
    \begin{cases} 0 & (x_i = 0),\\ 1 + \log_{10}|x_i| & (x_i \neq 0) \end{cases} \in \mathbb{R}_{\geq 0}$
  \item \textbf{Sign}: $s_i \in \{+, -, 0\}$ (3-class label)
  \item \textbf{Residues}: $r_i^{(m)} = x_i \bmod m$ for each $m \in \{2, 3, \ldots, 101\}$
        (100 independent classification targets)
\end{enumerate}
This decomposition separates magnitude, sign, and periodic arithmetic structure into
complementary supervision signals.

\subsection{Dual-Stream Representation}
\label{sec:dualstream}

%   f_i^mag = [v_i, 1[x_i>0], 1[x_i<0], 1[x_i=0]]
Two feature vectors are computed for each element $x_i$ prior to learnable embeddings.

\textbf{Magnitude features} $\mathbf{f}_i^{\text{mag}} \in \mathbb{R}^4$:
\[
  \mathbf{f}_i^{\text{mag}} = \bigl[v_i,\;\mathbf{1}[x_i > 0],\;\mathbf{1}[x_i < 0],\;\mathbf{1}[x_i = 0]\bigr].
\]
The last three components are a one-hot sign representation. For astronomically large integers
exceeding the \texttt{float64} range, $|x_i|$ is replaced by its decimal digit count.

\textbf{Modulo features} $\mathbf{f}_i^{\text{mod}} \in \mathbb{R}^{200}$:
For each modulus $m \in \{2, \ldots, 101\}$, let $r = x_i \bmod m$, and embed the residue as
a point on the unit circle:
\[
  \phi_m(r) = \left[\sin\!\left(\tfrac{2\pi r}{m}\right),\;\cos\!\left(\tfrac{2\pi r}{m}\right)\right] \in \mathbb{R}^2.
\]
Concatenating all 100 moduli gives $\mathbf{f}_i^{\text{mod}} \in \mathbb{R}^{200}$.
This sin/cos embedding is equivariant to the group structure of $\mathbb{Z}/m\mathbb{Z}$ and
avoids discontinuities at wrap-around boundaries.

The two feature vectors are projected to the model hidden dimension $d$ by independent
projection layers (two-layer MLP for magnitude; affine for modulo):
\[
  \mathbf{h}_i^{\text{mag}} = \mathrm{MLP}_{\text{mag}}\!\left(\mathbf{f}_i^{\text{mag}}\right),
  \quad
  \mathbf{h}_i^{\text{mod}} = W_{\text{mod}}\,\mathbf{f}_i^{\text{mod}} + \mathbf{b}_{\text{mod}},
  \quad \mathbf{h}_i^{\text{mag}},\,\mathbf{h}_i^{\text{mod}} \in \mathbb{R}^d.
\]

The two streams are fused by FiLM~\cite{perez2018film}: the modulo embedding generates
element-wise scale $\boldsymbol{\gamma}_i$ and shift $\boldsymbol{\beta}_i$:
\[
  \boldsymbol{\gamma}_i = W_\gamma\,\mathbf{h}_i^{\text{mod}},
  \quad
  \boldsymbol{\beta}_i = W_\beta\,\mathbf{h}_i^{\text{mod}},
  \quad
  \mathbf{e}_i = (1 + \boldsymbol{\gamma}_i) \odot \mathbf{h}_i^{\text{mag}} + \boldsymbol{\beta}_i.
\]
A ReLU is applied after the modulo projection, and dropout is inserted before FiLM.
Standard sinusoidal positional encodings are added to $\mathbf{e}_i$ before the encoder.

The overall architecture is illustrated in Fig.~\ref{fig:architecture}.

\begin{figure}[htbp]
  \centering
  \includegraphics[width=\linewidth]{figures/fig1_architecture}
  \caption{IntSeqBERT architecture. The Dual-Stream Embedding block projects
           $\mathbf{f}_i^{\mathrm{mag}}\!\in\!\mathbb{R}^4$ and
           $\mathbf{f}_i^{\mathrm{mod}}\!\in\!\mathbb{R}^{200}$ to $\mathbb{R}^d$ via
           $\mathrm{MLP}_{\mathrm{mag}}$ and $W_{\mathrm{mod}}$, then fuses them with FiLM:
           $\mathbf{e}_i = (1+\boldsymbol{\gamma}_i)\odot\mathbf{h}_i^{\mathrm{mag}}+\boldsymbol{\beta}_i$.
           Positional encodings are added before the Pre-LN Transformer encoder.
           Three prediction heads produce
           $\hat{v},\log\hat{\sigma}^2\!\in\!\mathbb{R}$ (magnitude regression),
           $\hat{s}\!\in\!\{+,-,0\}$ (sign classification), and
           $\hat{r}^{(m)}\!\in\!\{0,\ldots,m{-}1\}$ ($100\times$ modulo classification).}
  \label{fig:architecture}
\end{figure}

\subsection{Encoder, Heads, and Training}
\label{sec:encoder}

% Large (12L, d=768, 12h, 91.5M)
The fused sequence $(\mathbf{e}_1, \ldots, \mathbf{e}_L)$ is processed by a standard
Transformer encoder~\cite{vaswani2017attention} with Pre-Layer Normalisation~\cite{xiong2020layer}.
Three model sizes are evaluated (Table~\ref{tab:models}).

\begin{table}[htbp]
\caption{Model configurations.}
\label{tab:models}
\centering
\begin{tabular}{lcccr}
\hline
Config  & Layers & $d$ & Heads & Parameters \\
\hline
Small   & 6      & 256 & 4     & 6.4M       \\
Middle  & 8      & 512 & 8     & 29.0M      \\
Large   & 12     & 768 & 12    & 91.5M      \\
\hline
\end{tabular}
\end{table}

Let $\mathbf{z}_i \in \mathbb{R}^d$ be the encoder output at masked position $i$.
Three prediction heads are applied: (i)~\textbf{Magnitude head} --- a two-layer MLP
($d \to d \to 2$, ReLU) outputting $(\mu_i, \log\sigma_i^2)$, predicting $\hat{v}_i = \mu_i$
(a log-scale estimate of $v_i$; $\sigma_i^2 = \exp(\log\sigma_i^2)$ is used by the Solver);
(ii)~\textbf{Sign head} --- $\hat{s}_i = \operatorname{softmax}(W_{\text{sign}}\mathbf{z}_i)$;
(iii)~\textbf{Modulo head} --- 100 independent linear classifiers, one per modulus, with total
output dimension $\sum_{m=2}^{101}m = 5{,}150$.

The multi-task loss is $\mathcal{L} = \mathcal{L}_{\text{mag}} + \mathcal{L}_{\text{sign}} + 2\,\mathcal{L}_{\text{mod}}$.
The weight~2 on $\mathcal{L}_{\text{mod}}$ is intentionally set higher to emphasise the
newly introduced modulo stream; adaptive weighting was attempted but caused training instability.
$\mathcal{L}_{\text{mag}}$ is Huber loss; $\mathcal{L}_{\text{sign}}$ is cross-entropy over
three sign classes; $\mathcal{L}_{\text{mod}} = \frac{1}{100}\sum_{m=2}^{101}
\frac{1}{\ln m}\mathcal{L}_{\text{CE}}^{(m)}$, where $\mathcal{L}_{\text{CE}}^{(m)}$ is the
cross-entropy loss of the $m$-class classifier for modulus $m$, normalised by $\ln m$
(the maximum entropy of a uniform distribution over $m$ classes) to compensate for varying class counts.
All losses are computed only at masked positions.

\subsection{Baselines}
\label{sec:baselines}

\textbf{Vanilla Transformer} maps each integer to a token ID in a vocabulary of 20,003 entries
(non-negative values 0--19{,}999 plus \texttt{PAD}, \texttt{MASK}, \texttt{UNK}); negative
integers and values $\geq 20{,}000$ are replaced by \texttt{UNK}. The same three prediction heads are applied to the token
embeddings. The vocabulary size 20,003 was chosen so that memory consumption matches
IntSeqBERT under the 8\,GB VRAM constraint. The prior benchmark FACT~\cite{zurich-fact}
handles values up to several million, presupposing larger compute resources than our setting.

\textbf{Ablation (magnitude-only)} is identical to IntSeqBERT but uses only the magnitude
stream; the FiLM module is removed and $\mathbf{e}_i = \mathbf{h}_i^{\text{mag}}$. This
isolates the contribution of the modulo stream.

\subsection{Solver}
\label{sec:solver}

The pre-trained model outputs magnitude $(\mu, \log \sigma^2)$, sign, and modulo
distributions at masked positions. \textbf{Solver} recovers concrete integers from these
predictions.

The Solver derives the $3\sigma$ interval $[n_{\min}, n_{\max}]$ from the magnitude prediction
and dynamically selects one of three modes based on the search width
$\Delta n = |n_{\max} - n_{\min}|$ (Table~\ref{tab:solver}).

\begin{table}[htbp]
\caption{Solver operating modes.}
\label{tab:solver}
\centering
\begin{tabular}{llp{5.5cm}}
\hline
Mode     & Condition                              & Method \\
\hline
Dense    & $\Delta n \leq 10^6$                   & Enumerate all integers in range \\
Sieve    & $10^6 < \Delta n \leq 10^{14}$         & CRT beam search anchored on high-confidence moduli \\
CRT      & $\Delta n > 10^{14}$                   & Sparse CRT beam search to generate large integers directly \\
\hline
\end{tabular}
\end{table}

When the sign prediction is zero, the Solver immediately returns 0. If no valid candidate is
found in the search range, the call is recorded as \texttt{none}.

Each candidate $n$ is scored by $\text{score}(n) = \alpha_{\text{mag}} + 0.3\,\alpha_{\text{mod}}$,
where $\alpha_{\text{mag}} = -(v_n-\mu_i)^2/(2\sigma_i^2)$ and
$\alpha_{\text{mod}} = \sum_{m}\ln P(n\bmod m)$.
The coefficient 0.3 approximates the proportion of primes between 2 and 101
($26/100 = 0.26$), compensating for the information overlap among composite moduli
that share prime factors.
The top $k$ candidates are returned and evaluated as Solver Top-$k$ (Section~\ref{sec:solver-eval}).
